\section{Androids, Avatars, and Androphiles}

Most people think in sentence fragments. The more educated can think in complete sentences, or more rarely, in paragraphs.

\paragraph{Winding Down}
As we are getting close to the ¾ mark of 1001 posts, I have to think about wrapping things up. Just before \#1001, I will reveal the secret meaning of one of Rene Guenon's major themes. It should be obvious given everything he has written and he did provide some oblique clues. Gornahoor has actually provided some clues. Esoteric writings usually leave something out and Guenon leaves out what anyone familiar with his path should expect to see. There can be no hopelessness in the task of recovering tradition in the West, since hope is a virtue and the East is not bastion of tradition.

The three motive forces in the human mind are eros, thumos, and nous. The first two have a positive character when dominated by nous, but usually are allowed to run amok without the guidance of a higher intelligence. As spiritual training, readers should strive to determine which of these forces are being appealed to. Eros seems straight forward, but we are intending something more than simple appeals to prurient interest. Desire takes many forms. For example, in my fb feed, I often see links to sites on meditation that promise to lower blood pressure, to reduce anxiety, or, in general, to ``improve one's life". Obviously, that shows absolutely no understanding of the purpose of meditation and reveals the danger of learning meditation on one's own without a suitable guide. This is the viewpoint of the shudra, whose understanding is limited to the material and the subjective.

Perhaps the goal is rather to transcend such desires, to observe how they try to dominate consciousness, and then to reveal their emptiness.

In the areas we want to focus on here, thumos is appealed to in politically motivated writings. They are written in a way to appeal to anger and fear, to ``pump up". To do so, their technique is to attack the other side, however defined, while devoid of anything of positive content. Do we really need another exposé of the ``modern world", especially since almost everyone in the West is in love with it?

We, on the other hand, endeavor to appeal to the nous; that is why we don't ``prove" things empirically or logically (although they are useful in refutation). It is more important to ``see". Since the world of appearances in the reflection of something transcendent, we point to the transcendent. While others work themselves up over the shadow play on the wall, we want readers to see the puppets rather than their shadows. Therefore we leave you with two homework tasks:

1) When reading something, or better yet if you can maintain the requisite detachment and concentration, when engaging with someone in person, notice which force is dominant in them.

2) Similarly, try to determine which caste a person probably should belong to. Schuon's schema of objective/subjective and idealist/materialist is a good way to start.

\paragraph{Satanism, the Vatican, and the Old Religion}
There was a report recently about a group of homosexual priests in Rome who maintained a ring of catamites. To juice it up, the story claimed that this group provided consecrated hosts to Satanists. The Italian police quickly hushed it up as the ravings of a bitter ex-priest. Nevertheless, there is a positive element in that someone actually believes in the real presence of Christ in the Eucharist, showing that Satan can't be fooled. Beyond that, it sounds like some new Dan Brown novel about androphiles following the left hand path; there are enough of them in the counter tradition.

Actually, Satan's religion is the first religion, pre-existing the primordial revelation given to Adam; that is why the witches call it the ``Old Religion". Its creed is simple: deny God and deny God's cosmic law. These denials take several forms, from simple atheism to deliberate and conscious rejection. The Old Religion sees matter as the prison of man, rather than the symbol of transcendence. Adam's revelation understands the problem as arising from the human will, the rejection of the law, which is defined as ``sin", i.e., deviating from that cosmic norm.

For example, I recently came across a discussion of some recent SCOTUS decisions; the author decided that some things can be ``natural" but not ``normative" and believed that the observation of forest animals is a sufficient guide to understand human behavior. This presumes, falsely, that nature is independent and can be fully understood apart from creation. The author claims that the idea of ``sin" clouds the issue, and I agree that the term is overworked and not always understood, yet it is indispensable. For the Old Religionist, no appeal to any law, not even that of forest animals, is persuasive.

\paragraph{Avatars and Immortality}
Recently, there was a conference in New York called the 2045 Initiative. Their stated goal is to create artificial humans, or androids, and then ``upload" the contents of a human mind into that being. Of course, there is absolutely no reason to believe that such a project is even feasible and the participants provided none. They simply assume a Cartesian view of things.

I came across one writer who ``defined" consciousness as the desires, emotions, images, thoughts, in short, all the contents in the mind. Unfortunately, as consciousness is not an object in the world, no such definition is possible. Consciousness is transcendent, even to such ``mental" contents, and is always the subject. Does everyone reading this understand that notion? Ask for clarification, if not.

There is a more serious difficulty in that significant parts of the contents of the soul actually belong to the body; hence, they cannot be separated from the body. Perhaps logical forms of thought can be, but that is not the human person, but rather a simulation. Would they try to eliminate all irrational desires, urges, unsupportable opinions, prejudices, etc.? For most men, there would not be much left over.

What they are attempting is to create a copy of a copy. Nature is God's creation and man is part of that nature. The android is an artifact, something not natural in the same way. Specifically, there is no divine idea of it, but that is a complete topic in itself.

\paragraph{Drunk Lamas and Pusillanimity}
Since a Lama is the incarnation of an aspect of the Buddha, then accidental qualities such as his dipsomania cannot affect that. There is a similar teaching in the Western tradition that is called ``ex opere operato"; the rites of a priest are effective regardless of his personal failings. But the more interesting metaphysical question, however, is why was such an imperfect being chosen as the lama? On a more elevated level, we can ask the related question: ``why was the woman St. Joan of Arc chosen to save France?"

Well, we know that there is a revealed truth that explains it: ``Many are called but few are chosen." Specifically, many men, of better character than the Mongolian lama, were called to fulfill that role. Yet they passed the cup, assuming George would do it\footnote{\url{https://www.oldtimeradiodownloads.com/crime/let-george-do-it}}, until finally one spirit accepted.

So we can assume that many men were offered the chance to ``save" France until there were no men left. Pusillanimity is the vice that holds a man back from actualizing all his possibilities. Often, he will justify it by the claim to be ``humble" or offer some other excuse. The call often arises as a whisper out of the Silence. If your mind is so cluttered with loud, persistent, meaningless and strength-sapping thoughts then that whisper will be missed. Actually the pusillanimous man hopes it will be missed.

\paragraph{The Boy who Cried Wolf}
I watched the 1974 version of the Great Gatsby last week and was struck by the allusion to Lothrop Stoddard; apparently, that made no impact on me when I first saw it or read the book. Considering that his major opus, and those of his spiritual father Madison Grant, were written nearly a century ago, it is interesting to compare similar ideas now. Then Stoddard and Grant were respectable intellectuals, and the latter was even chummy with some presidents. Nowadays, however, neither one of them could even get a cookbook published.

If circumstances were so auspicious at that time, then what chance do similar ideas have of taking root now? There is none, since there is no common identity and this is proved by all the appeals to identity. This is actually an interesting topic when viewed from the right and deserves a post dedicated to it. It will be called: Order, Authority, and Unity.



\flrightit{Posted on 2013-07-07 by Cologero }
